\documentclass[11pt]{article}
\usepackage[a4paper,margin=26mm]{geometry}
\usepackage{amsmath,amssymb,amsthm,hyperref,microtype,booktabs}
\hypersetup{hidelinks}
\newtheorem{theorem}{Theorem}
\newtheorem{lemma}[theorem]{Lemma}
\newtheorem{corollary}[theorem]{Corollary}
\newtheorem{proposition}[theorem]{Proposition}
\newcommand{\C}{\mathbb C}
\newcommand{\R}{\mathbb R}
\newcommand{\T}{\mathbb T}
\newcommand{\Ima}{\operatorname{Im}}
\newcommand{\rank}{\operatorname{rank}}
\newcommand{\diag}{\operatorname{diag}}
\title{Odd Jordan blocks determine the imaginary rank of a limiting matrix Riccati solution}
\author{}
\date{}
\begin{document}
\maketitle
\begin{abstract}
Let $X_\eta$ be the stabilizing, positive-imaginary-part solution of
$X_\eta+A^\top X_\eta^{-1}A=Q+i\eta I$, where $A$ is real, $Q$ is real symmetric, and $\eta>0$.
Assume that $X_\eta$ has a finite invertible limit $X$.
We prove that the rank of $\Ima X$ is one half the number of odd-sized Jordan blocks of
$\lambda^2A^\top-\lambda Q+A$ on the unit circle. Blocks, not their algebraic multiplicities, are counted.
The proof uses local Rellich eigenbranches, confluent interpolation of stable eigenvectors,
and a positive Cauchy kernel furnished by the Stein identity.
This gives the semisimple rank equality without an auxiliary regularity or non-splitting condition,
and treats defective unit-circle eigenvalues as well.
An exact three-by-three example has two Jordan blocks of size three and limiting imaginary rank one.
The finite-invertible-limit assumption is essential to the scope of the result and is not asserted to hold universally.
\end{abstract}

\section{Setting and main result}
Let $A\in\R^{n\times n}$ and $Q=Q^\top\in\R^{n\times n}$.
For $\eta>0$, write $Q_\eta=Q+i\eta I$ and let $X_\eta=X_\eta^\top$ satisfy
\begin{equation}\label{eq:riccati}
X_\eta+A^\top X_\eta^{-1}A=Q_\eta,
\qquad H_\eta=\Ima X_\eta>0,
\qquad \rho(S_\eta)<1,\quad S_\eta=X_\eta^{-1}A.
\end{equation}
Here $\Ima Y=(Y-Y^*)/(2i)$. For a complex symmetric matrix this is a real symmetric matrix.
Existence and uniqueness of this stabilizing solution for $\eta>0$ are established in
\cite[Theorem~2.2 and Remark~2.1]{GKL}. Throughout the main result we assume
\begin{equation}\label{eq:limit}
X=\lim_{\eta\downarrow0}X_\eta\quad\text{exists as a finite invertible matrix}.
\end{equation}
Let
\[
P_\eta(\lambda)=\lambda^2A^\top-\lambda Q_\eta+A,\qquad P=P_0.
\]
A Jordan block of the regular quadratic polynomial $P$ means a finite elementary divisor,
or, equivalently, the corresponding Jordan block in a strong linearization.
Its size is the positive partial multiplicity at that eigenvalue.

\begin{theorem}\label{thm:main}
Assume \eqref{eq:riccati}--\eqref{eq:limit}.
Let $b_{\rm odd}$ be the total number of odd-sized Jordan blocks of $P$ at eigenvalues
$\lambda$ with $|\lambda|=1$, counting each block once.
Then $b_{\rm odd}$ is even and
\begin{equation}\label{eq:rank}
\rank(\Ima X)=\frac{b_{\rm odd}}2.
\end{equation}
The same rank formula holds for $\Ima(X^{-1})$.
In particular, $\Ima X=0$ if and only if all unit-circle Jordan blocks have even size.
\end{theorem}

Guo--Kuo--Lin \cite[Theorem~3.3]{GKL} prove that, if $2m$ is the total algebraic
multiplicity on the unit circle, then $\rank(\Ima X)\le m$.
They obtain equality under semisimplicity and an $O(\eta)$ regularity assumption,
or under semisimplicity and a non-splitting condition on each repeated eigenvalue cluster.
Theorem~\ref{thm:main} gives a finer blockwise count. In the semisimple case every block
has size one, so it gives equality $\rank(\Ima X)=m$ directly.

The algebraic perturbation facts used below are local. No globally single-valued Rellich
eigenbasis on the entire unit circle is assumed. Local analytic Hermitian eigendecompositions
and the identification of their zero orders with partial multiplicities are standard;
see \cite[Theorem~1.1 and Section~5]{BN}.

\section{The Stein identity and the stable graph}
Taking imaginary parts in \eqref{eq:riccati}, using
$\Ima(X_\eta^{-1})=-X_\eta^{-*}H_\eta X_\eta^{-1}$, gives
\begin{equation}\label{eq:Stein}
H_\eta-S_\eta^*H_\eta S_\eta=\eta I.
\end{equation}
Also, symmetry of $X_\eta$ gives the factorization
\begin{equation}\label{eq:factor}
P_\eta(\lambda)=(I-\lambda S_\eta^\top)X_\eta(S_\eta-\lambda I).
\end{equation}
The same factorization holds at $\eta=0$, with $S=X^{-1}A$.
It shows that $P$ is regular; the two determinant factors are nonzero polynomials,
and $X$ is invertible. It also shows that $P_\eta$ has exactly $n$ stable eigenvalues,
counted with multiplicity, and the other $n$ eigenvalues are outside the unit disk
on the extended plane. Zero and infinite eigenvalues are allowed.

Use the linearization
\[
M_\eta=\begin{pmatrix}A&0\\Q_\eta&-I\end{pmatrix},
\qquad L=\begin{pmatrix}0&I\\A^\top&0\end{pmatrix}.
\]
Its stable right deflating subspace is
\begin{equation}\label{eq:graph}
\mathcal W_\eta=\operatorname{range}\begin{pmatrix}I\\X_\eta\end{pmatrix},
\end{equation}
because
\[
M_\eta\begin{pmatrix}I\\X_\eta\end{pmatrix}
=L\begin{pmatrix}I\\X_\eta\end{pmatrix}S_\eta.
\]
The columns on the right remain independent after multiplication by $L$, since the upper
block of $L[I;X_\eta]$ is $X_\eta$.
A finite eigenpair $P_\eta(\lambda)v=0$ lifts to the pencil eigenvector
\begin{equation}\label{eq:lift}
\zeta_\eta(\lambda,v)=\begin{pmatrix}v\\(Q_\eta-\lambda A^\top)v\end{pmatrix}.
\end{equation}
For a stable eigenvalue, its upper block $v$ is therefore a right eigenvector of $S_\eta$.
If $S_\eta v_a=\lambda_a v_a$ and $S_\eta v_b=\lambda_b v_b$, then
\eqref{eq:Stein} yields the exact identity
\begin{equation}\label{eq:kernel}
v_a^*H_\eta v_b=\frac{\eta\,v_a^*v_b}{1-\overline{\lambda_a}\lambda_b}.
\end{equation}

\section{Local branches and the number of stable roots}
Fix a unit-circle eigenvalue $\lambda_0=e^{i\theta_0}$ and put
\[
T(t)=e^{i(\theta_0+t)}A^\top+e^{-i(\theta_0+t)}A-Q.
\]
For real $t$, $T(t)$ is Hermitian. By the local Rellich theorem,
\begin{equation}\label{eq:Rellich}
T(t)=U(t)\diag(d_1(t),\ldots,d_n(t))U(t)^{-1},
\qquad U(\overline t)^*U(t)=I,
\end{equation}
where the functions extend holomorphically to a complex neighborhood of zero, and each
$d_j$ is real on the real line. The second identity follows by analytic continuation
from real $t$; $U(\overline t)^*$ denotes the holomorphic reflection of $U$.
For a branch vanishing at zero, write
\begin{equation}\label{eq:branch}
d_j(t)=a_jt^{\ell_j}+O(t^{\ell_j+1}),\qquad a_j\in\R\setminus\{0\},\quad \ell_j\ge1.
\end{equation}
Regularity excludes an identically zero branch. The orders $\ell_j$ are exactly the
positive partial multiplicities of $P$ at $\lambda_0$: multiplication by the nonvanishing
scalar $e^{i(\theta_0+t)}$ and the local analytic change from $t$ to $\lambda$ preserve
partial multiplicities, and \eqref{eq:Rellich} is an analytic equivalence to a diagonal matrix.

Importantly, the regularization is scalar. Thus the same $U(t)$ diagonalizes
$T(t)-i\eta I$, and the perturbed roots from branch $j$ solve $d_j(t)=i\eta$.
Put $\varepsilon=\eta^{1/\ell_j}$. The roots have the form
\begin{equation}\label{eq:rootexp}
t_a=\varepsilon z_a+O(\varepsilon^2),\qquad a_jz_a^{\ell_j}=i.
\end{equation}
This follows by the analytic implicit function theorem after substituting $t=\varepsilon z$;
the limiting roots $z_a$ are distinct and nonzero. All the resulting roots are simple
within this branch for sufficiently small positive $\eta$.
The corresponding $\lambda=e^{i(\theta_0+t)}$ is stable exactly when $\Ima t>0$.
No $z_a$ is real. Counting the upper-half-plane roots of $a_jz^{\ell_j}=i$ gives
\begin{equation}\label{eq:r}
r_j=\begin{cases}
\ell_j/2,&\ell_j\text{ even},\\
(\ell_j+1)/2,&\ell_j\text{ odd and }a_j>0,\\
(\ell_j-1)/2,&\ell_j\text{ odd and }a_j<0.
\end{cases}
\end{equation}
Coincidences of roots belonging to different branches do not obstruct the construction:
at a common root the corresponding columns of $U(t)$ are independent.

\section{Confluent stable bases}
We record explicitly the interpolation fact that controls the degenerating eigenvector bases.

\begin{lemma}\label{lem:interpolation}
Let $t_1,\ldots,t_r$ be distinct nodes with $t_a=\varepsilon z_a+O(\varepsilon^2)$,
where the $z_a$ are distinct and $\varepsilon\downarrow0$.
Set $V_\eta=(t_a^k)_{a=1,\ldots,r;\,k=0,\ldots,r-1}$.
If $f(t)=\sum_{k\ge0}f_kt^k$ is a vector-valued holomorphic function, then
\[
[f(t_1),\ldots,f(t_r)]V_\eta^{-\top}\longrightarrow[f_0,\ldots,f_{r-1}].
\]
For a function holomorphic in two variables near $(0,0)$, applying these interpolation
transformations in both variables gives a bounded coefficient matrix as both sets of
nodes approach zero. The two sets may approach at different rates.
\end{lemma}
\begin{proof}
The displayed columns are the coefficients of the degree-$r-1$ polynomial interpolating
$f$ at the nodes. For a small fixed contour enclosing all nodes, the divided-difference
formula bounds these coefficients uniformly and shows that they tend to the Taylor
coefficients as the nodes coalesce. More explicitly, the Newton coefficients are contour
integrals of $f(z)/\prod_{a=1}^{k+1}(z-t_a)$; they tend to $f_k$, and conversion from
Newton to monomial coefficients tends to the identity. Applying the same contour formula
successively in two variables proves the last assertion, uniformly and without a condition
on relative rates. The result also holds for a uniformly holomorphic family converging
to $f$, as used for the lifts below.
\end{proof}

For a branch $j$ with $r_j>0$, let $v_j(t)=U(t)e_j$, restrict \eqref{eq:rootexp} to its
$r_j$ stable roots, and define
\begin{align*}
E_{j,\eta}&=[v_j(t_1),\ldots,v_j(t_{r_j})],\\
B_{j,\eta}&=E_{j,\eta}V_{j,\eta}^{-\top},\qquad
V_{j,\eta}=(t_a^k)_{a,k}.
\end{align*}
Apply the same transformation to the lifted vectors
\[
\zeta_{j,\eta}(t)=
\begin{pmatrix}v_j(t)\\(Q+i\eta I-e^{i(\theta_0+t)}A^\top)v_j(t)\end{pmatrix}.
\]
The resulting $2n\times r_j$ matrix tends to the first $r_j$ Taylor coefficients of
$\zeta_{j,0}(t)$. These coefficients form, up to an invertible triangular change of
basis induced by $\lambda=e^{i(\theta_0+t)}$, the first $r_j$ vectors of the Jordan chain
associated with this partial multiplicity. Indeed,
\[
(M_0-e^{i(\theta_0+t)}L)\zeta_{j,0}(t)
=\begin{pmatrix}e^{i(\theta_0+t)}d_j(t)v_j(t)\\0\end{pmatrix}
=O(t^{\ell_j}).
\]
The analytic equivalence \eqref{eq:Rellich} supplies a complete canonical system of root
functions, so the full lifted chains are independent; any collection of their prefixes
is independent as well. This is also the usual chain construction from the local Smith
form of a regular matrix polynomial. The $i\eta I$ term in the lift contributes $i\eta$
times the interpolated upper block, which tends to zero.

The part of the pencil spectrum strictly inside the unit disk is separated from the
boundary. Its full right deflating subspace has an analytic basis $W_{s,\eta}$ near zero,
even when individual eigenvalues are defective. For example, a fixed shift off the spectrum
reduces this assertion to an analytic Riesz projection of $(M_\eta-\sigma L)^{-1}L$.
Adjoin to this basis all the normalized boundary blocks just constructed. The result is
\begin{equation}\label{eq:W}
W_\eta=\begin{pmatrix}B_\eta\\C_\eta\end{pmatrix}\longrightarrow
W_0=\begin{pmatrix}B_0\\C_0\end{pmatrix},
\qquad \rank W_0=n,\qquad \operatorname{range}W_\eta=\mathcal W_\eta.
\end{equation}
The rank follows from independence of the chain prefixes and direct sums between distinct
pencil eigenvalues. There are $n$ columns because the stable dimension is $n$ for $\eta>0$.
The graph property gives $C_\eta=X_\eta B_\eta$ and $B_\eta$ invertible. Taking limits yields
$C_0=XB_0$. If $B_0u=0$ then $W_0u=0$, so $u=0$. Consequently
\begin{equation}\label{eq:Binv}
B_0\text{ is invertible}.
\end{equation}
This is the precise place where the finite limit in \eqref{eq:limit} is used.

\section{The limiting Gram matrix}
We compute $B_\eta^*H_\eta B_\eta$ in the basis of \eqref{eq:W}.
The strictly stable block contributes zero in the limit. In fact its upper block satisfies
$S_\eta B_{s,\eta}=B_{s,\eta}R_{s,\eta}$, with $R_{s,\eta}$ tending to a matrix with spectral
radius strictly less than one. Restricting \eqref{eq:Stein} to this block gives
\[
G_{s,\eta}-R_{s,\eta}^*G_{s,\eta}R_{s,\eta}
=\eta B_{s,\eta}^*B_{s,\eta},\qquad G_{s,\eta}=B_{s,\eta}^*H_\eta B_{s,\eta}.
\]
The inverse Stein operator is uniformly bounded near the limit, hence $G_{s,\eta}=O(\eta)$.
All mixed blocks with this block tend to zero by positive semidefiniteness and the
Cauchy--Schwarz inequality; the other blocks remain bounded by \eqref{eq:limit} and
\eqref{eq:W}.

Next consider two branches at the same $\theta_0$. Define the holomorphic row function
$u_j(s)=v_j(\overline s)^*$. Since $u_i(t)v_j(t)=\delta_{ij}$, division by $t-s$ gives
\begin{equation}\label{eq:numerator}
u_i(s)v_j(t)=\delta_{ij}+(t-s)F_{ij}(s,t)
\end{equation}
for a holomorphic function $F_{ij}$. The denominator in \eqref{eq:kernel} is
$1-e^{i(t-s)}$. It has a simple zero at $t=s$.
Therefore, for different branches $i\ne j$, the entire Gram kernel is $\eta$ times a
holomorphic function of $(s,t)$. Lemma~\ref{lem:interpolation} shows that its mixed block
in the normalized basis is $O(\eta)$, even if the branches have different orders.
For different boundary points the denominator is already nonzero near the two centers,
so the same conclusion holds. For one branch, \eqref{eq:numerator} shows that its Gram
block equals the transformed universal kernel
\begin{equation}\label{eq:universal}
K_{ab}=\frac{\eta}{1-e^{i(t_b-\overline{t_a})}}
\end{equation}
plus a block of order $O(\eta)$.

It remains to calculate the rank contribution from \eqref{eq:universal}.
Drop the branch index, put $\ell=\ell_j$, $r=r_j$, and suppose $r>0$.
With $\varepsilon=\eta^{1/\ell}$ and $t_a=\varepsilon z_a+O(\varepsilon^2)$,
\begin{equation}\label{eq:Ckernel}
K=\varepsilon^{\ell-1}(C+O(\varepsilon)),
\qquad C_{ab}=\frac{i}{z_b-\overline{z_a}}.
\end{equation}
The matrix $C$ is positive definite: it is the Gram matrix of the linearly independent
functions $e^{iz_at}$ in $L^2(0,\infty)$, since all $\Ima z_a>0$ and
\[
\int_0^\infty e^{-i\overline{z_a}t}e^{iz_bt}\,dt
=\frac{i}{z_b-\overline{z_a}}.
\]
Write
\[
V_\eta=\widehat V_\eta D_\varepsilon,
\qquad (\widehat V_\eta)_{ak}=(t_a/\varepsilon)^k,
\qquad D_\varepsilon=\diag(1,\varepsilon,\ldots,\varepsilon^{r-1}).
\]
Then $\widehat V_\eta$ tends to the invertible Vandermonde matrix $\widehat V=(z_a^k)$.
The Gram block after the normalization $E_\eta V_\eta^{-\top}$ is
\begin{align}
\overline{V_\eta}^{-1}KV_\eta^{-\top}
&=\varepsilon^{\ell-1}D_\varepsilon^{-1}
   (D+O(\varepsilon))D_\varepsilon^{-1},\label{eq:scaledGram}\\
D&=\overline{\widehat V}^{-1}C\widehat V^{-\top}>0.\notag
\end{align}
For indices $k,l=0,\ldots,r-1$, its $(k,l)$ entry is
$\varepsilon^{\ell-1-k-l}(D_{kl}+O(\varepsilon))$.
By \eqref{eq:r}, the smallest exponent is
\[
\ell+1-2r=
\begin{cases}
1,&\ell\text{ even},\\
0,&\ell\text{ odd and }a_j>0,\\
2,&\ell\text{ odd and }a_j<0.
\end{cases}
\]
Hence an even branch or a negative odd branch has zero limiting block.
A positive odd branch has exactly one surviving entry, the last diagonal entry
$D_{r-1,r-1}>0$, and contributes rank one. When $r=0$ there is no block and no contribution.
All remaining mixed blocks tend to zero, as already shown. Together with
\eqref{eq:Binv}, this proves the intermediate formula
\begin{equation}\label{eq:signedrank}
\rank(\Ima X)=\#\{\text{unit-circle Rellich zero branches with odd order and }a_j>0\}.
\end{equation}

\section{Real symmetry and the half-count}
For the global Hermitian trigonometric polynomial
$\mathcal T(\theta)=e^{i\theta}A^\top+e^{-i\theta}A-Q$, reality gives
\[
\mathcal T(-\theta)=\overline{\mathcal T(\theta)}=\mathcal T(\theta)^\top.
\]
Thus its multiset of analytic eigenvalue germs at $-\theta_0$ is obtained from that
at $\theta_0$ by $t\mapsto-t$. A branch of order $\ell$ and leading coefficient $a$
is paired with a branch of order $\ell$ and coefficient $(-1)^\ell a$.
At distinct conjugate unit-circle points this pairs the two collections of germs.
At $\theta_0=0$ or $\pi$ the same statement holds within the collection at that point,
using $2\pi$-periodicity. Odd-order germs cannot be fixed by this operation because their
leading coefficient changes sign; their positive and negative counts are equal.
Consequently the total number of positive odd branches is exactly $b_{\rm odd}/2$.
Equation~\eqref{eq:signedrank} proves \eqref{eq:rank} and evenness of $b_{\rm odd}$.
Finally
\[
\Ima(X^{-1})=-X^{-*}(\Ima X)X^{-1}
\]
is an invertible congruence up to sign, so its rank is the same. This completes the proof.

\section{Semisimple regularity as a corollary}
\begin{corollary}\label{cor:semisimple}
Under the assumptions of Theorem~\ref{thm:main}, suppose all unit-circle eigenvalues are
semisimple, of total algebraic multiplicity $2m$. Then
\[
\rank(\Ima X)=m,\qquad X_\eta-X=O(\eta),\qquad S_\eta-X^{-1}A=O(\eta).
\]
\end{corollary}
\begin{proof}
The rank conclusion follows from Theorem~\ref{thm:main}.
Every vanishing branch now has $\ell=1$ and $a_j\ne0$.
Its stable root, when it exists, is analytic in $\eta$ and is $O(\eta)$ from its limit,
as is its lifted eigenvector. Thus the boundary blocks and the separated strictly stable
block give $W_\eta=W_0+O(\eta)$ in \eqref{eq:W}.
The invertibility of $B_0$ then yields
$X_\eta=C_\eta B_\eta^{-1}=C_0B_0^{-1}+O(\eta)=X+O(\eta)$.
The final assertion follows by inversion at the invertible matrix $X$.
\end{proof}
This proves directly that the auxiliary regularity assumption in
\cite[Theorem~3.3(ii)]{GKL} is automatic under the other stated hypotheses.
It also permits a repeated semisimple cluster to split in both directions.

\section{An exact defective example}
Consider
\begin{equation}\label{eq:exampleAQ}
A=\begin{pmatrix}-1&0&-1/2\\0&0&-1/2\\1/2&-1/2&0\end{pmatrix},
\qquad
Q=\begin{pmatrix}-2&0&0\\0&0&-1\\0&-1&-1\end{pmatrix}.
\end{equation}
Writing $c=\cos\theta$ and $s=\sin\theta$ gives
\[
\mathcal T(\theta)=
\begin{pmatrix}2(1-c)&0&i s\\0&0&1-c\\-i s&1-c&1\end{pmatrix},
\qquad
\det P(\lambda)=\frac{(\lambda-1)^6}{4}.
\]
The two vanishing eigenbranches at zero have leading terms $\pm\theta^3/2$.
To verify the orders and signs, eliminate the last row and column at eigenvalue zero;
the Schur complement is
\[
\begin{pmatrix}(1-c)^2&-i s(1-c)\\i s(1-c)&-(1-c)^2\end{pmatrix}
=\frac{\theta^3}{2}\begin{pmatrix}0&-i\\i&0\end{pmatrix}+O(\theta^4).
\]
In the eigenvalue equation the additional mass factor tends to the identity, so these
are also the leading coefficients of the two small eigenvalues. Hence $P$ has exactly
two Jordan blocks of size three at $1$.

The stable solution has the finite limit
\begin{equation}\label{eq:exampleX}
X=\begin{pmatrix}-1&0&1/2\\0&0&-1/2\\1/2&-1/2&(-1+i)/2\end{pmatrix}.
\end{equation}
Here is a graph verification identifying the limit, rather than merely checking the
unregularized equation. The positive branch has $v_+(0)=(1,i,0)^\top$ up to normalization,
and its first derivative has third entry $i$, while its first two entries are immaterial
modulo the two vectors $v_+(0)$ and $v_-(0)=(1,-i,0)^\top$.
The positive branch contributes a chain prefix of length two and the negative branch
one of length one. Their upper blocks therefore span $\C^3$.
The normalized stable graph consequently has an invertible upper limit and determines
$X$ by its limiting lower block. On $e_1,e_2$ the resulting map is $Q-A^\top$; on $e_3$ it is
\[
(Q-A^\top)e_3-A^\top(e_1+i e_2)
=(1/2,-1/2,(-1+i)/2)^\top.
\]
This is precisely \eqref{eq:exampleX}.
Independently, direct exact calculation gives
\[
\det X=\det A=1/4,\quad
X+A^\top X^{-1}A=Q,\quad
X^{-1}A=\begin{pmatrix}1&0&1\\0&1&i\\0&0&1\end{pmatrix},\quad
\Ima X=\diag(0,0,1/2).
\]
Thus the limiting imaginary rank is one, not half the total unit-circle algebraic
multiplicity (which would be three). It agrees with half the number of odd blocks.
The exact identities and the Jordan structure are checked by
\path{code/nano_rank_certificate.py}; optional QZ experiments are reported separately
and are not used as a proof.

\section{Scope, reproducibility, and review}
The theorem is conditional on a finite invertible limiting solution. That assumption
cannot simply be dropped. For example,
\[
A=\begin{pmatrix}0&1\\0&0\end{pmatrix},\quad Q=0,\qquad
X_\eta=\begin{pmatrix}i\eta&0\\0&i(\eta+\eta^{-1})\end{pmatrix}
\]
satisfies \eqref{eq:riccati} with $\rho(S_\eta)=0$, but diverges as $\eta\downarrow0$.
Here $\det P_0(\lambda)=-\lambda^2$ and there are no unit-circle eigenvalues.
Neither semisimplicity nor absence of boundary eigenvalues alone guarantees the limit.

The scalar regularization $i\eta I$ is also used essentially: it preserves the local
Rellich eigenvectors and gives the scalar kernel in \eqref{eq:kernel}.
A different positive-definite regularization may be reduced by congruence, but no claim
about arbitrary singular or indefinite regularizations is made here.
The manuscript gives an analytic proof, supplemented by exact checks of its defective
example; a finite set of computations is not a verification of the general theorem.
The most delicate steps for independent review are the normalized Jordan-chain basis
and the two-variable interpolation argument for mixed blocks.
External proof review, priority review, and matching to the canonical repository wording
are still required before a catalog status change. This is an unreviewed solution claim,
not an externally validated resolution.

\begin{thebibliography}{9}
\bibitem{GKL}
C.-H. Guo, Y.-C. Kuo, and W.-W. Lin,
\emph{On a nonlinear matrix equation arising in nano research},
SIAM Journal on Matrix Analysis and Applications \textbf{33}(1) (2012), 235--262.
\href{https://doi.org/10.1137/100814706}{doi:10.1137/100814706}.
Author preprint:
\url{https://uregina.ca/~chguo/simax81470.pdf}.
\bibitem{BN}
G. Barbarino and V. Noferini,
\emph{On the Rellich eigendecomposition of para-Hermitian matrices and the sign
characteristics of $*$-palindromic matrix polynomials},
arXiv:2211.15539 (2022), Theorem~1.1 and Section~5.
\url{https://arxiv.org/abs/2211.15539}.
\end{thebibliography}
\end{document}
